% !TEX root = ../paper.tex
\section{Analysis of $\ts$ for arbitrary $3$-CNFs} \label{sec:non-monotone}



In this section, we analyze transversal trees for arbitrary $3$-CNFs and prove \cref{thm:main-result}. We will introduce few more ideas in addition to those introduced in \cref{sec:monotone}.

Throughout this section, we fix a $3$-CNF $F = (X, \cC)$ and let $T$ be its canonical transversal tree with root node $r$. Let the number of maximally disjoint width $3$ clauses used to develop $F$ be $m$. Let $X_D\subset X$ denote set of variables that appeared in this set of $m$ disjoint clauses. 
We also note that unlike in \cref{sec:monotone}, the clauses used to develop $T$ may not have width exactly $3$.


% \begin{definition}[Disjoint depth]
% For a \emph{canonical} transversal tree $T$, we define \emph{disjoint depth} to be the number of maximally disjoint clauses of width $k$ used to create $T$.
% \end{definition}

% Note that there can be other disjoint monotone clauses of width at most $k-1$ in $T$ but we insist they not be counted amongst these $m$ disjoint monotone clauses.

\subsection{Slight modification to canonical ordering and $\ts$}

We extend the conditions on the canonical ordering $\Pi$ of clauses and how $\ts$ uses $\Pi$. We order clauses in $\Pi$ so that all maximally disjoint width $3$ clauses appear first, followed by all width $3$ monotone clauses, followed by all other clauses.
For a node $u$ at level below $m$, we impose that instead of choosing the first unsatisfied monotone clause from $\Pi$, $u$ instead chooses an unsatisfied width $3$ monotone clause $C$ from $\Pi$ such that $C$ does not contain any variable $x\in X_D$ that has appeared twice in the shoot $S_{ru}$. If such a clause does not exist, then $C$ can pick the first unsatisfied monotone clause from $\Pi$.




\subsection{Fullness and Double marking}

We reuse the notion of weight here and observe that all basic facts and basic lemmas from monotone analysis apply here as well. We introduce one more definition related to this:

\begin{definition}[Uniform Weight]
Let $S = S_{uv}$ be a shoot in $T$ of length $\ell$. 
Let $a$ be the number of edges in the shoot $S_{uv}$.
The \emph{uniform weight} of $S$ denoted by $W^{+}(S) = W(S) + 3\ell - a$.
\end{definition}

We can similarly find a lower bound to this quantity for every root to leaf path:

\begin{fact}
\label{fact:shoot-uniform-extended-weight-lower-bound}
%Let $T$ be a depth-$d$ transversal tree for a monotone $3$-CNF. 
Every root to leaf shoot in $T$ has a uniform weight of
at least $3t - n$.
\end{fact}

\begin{proof}
Let $S$ be arbitrary root to leaf path with $a$ edges.
As there are only $n$ distinct edge labels, at least $a-n$ labels appear at least twice. So, $W(S) \ge a - n$.
Since the depth of $T$ is $t$, we infer that $W^{+}(S) = W^+(S) + 3t - a \ge 3t - n$ as desired.
\end{proof}

We extend the idea of markings and introduce double markings:

\begin{definition}[Double marking]
We say an edge $e\in T$ is doubly marked if $|M(e)| \ge 2$.
Let $P = P_{uv}$ be a path from $u$ to $v$. We write $W_{\ge 2}(P)$ to denote number of edges $e$ in $P$ such that $|M(e)| \ge 2$.
\end{definition}


Recall that when proving \cref{pess-ntran}, we took advantage of the fact that any monotone $3$-CNF $G$ with $\trn(G) = t$ will contain $\frac{t}{3}$ disjoint monotone clauses. However, there is no such guarantee for arbitrary $3$-CNFs. Observe that if the number of maximally disjoint monotone clauses is small, then many clauses of width $3$ will intersect with it and this will cause many edges to be doubly marked. We formalize this intution and introduce a new parameter called \emph{fullness} that will help keep track of this:

\begin{definition}[Fullness]
Let $u$ be arbitrary node at level $\ge m$ in $T$. Let $u_m$ be the node at level $m$ along the path $P_{ru}$. Then, \emph{fullness} of the shoot $S_{ru}$ is defined as
\[
\openfunction(S_{ru}) := |\{x\in X_D\setminus Q_{u_m}: \exists e = (a, b)\in S_{ru}, depth(a) \ge m, Q_e = x\}|,
\]
i.e., the number of variables in $X_D$ that are not along the path in the first $m$ levels and appear as labels of some edge in the shoot after level $m$. For arbitrary nodes $u, v$ where $u$ is ancestor of $v$ and $u$ appears at level $\ge m$, we define $\openfunction(S_{uv}) = \openfunction(S_{rv}) - \openfunction(S_{ru})$.
\end{definition}

This parameter is useful because every node after level $m$ will have at least one edge that will either `add to' fullness or at least one edge that will have marking set of size at least $2$. The edges that are `doubly marked' will contribute very little to the recursion. Moreover, we will see that fullness of any root to leaf shoot is at most $2m$. As $m$ is small, very few nodes will be such that they will not contain any doubly marked edges. This is our key insight and we fomrally prove this now.

\begin{fact}
\label{fact:disjoint-intersect}
   Let $M$ denote the set of all monotone clauses of width $3$ in $F$. Then, every clause $C\in M$ must contain at least one variable $x$ such that $x\in X_D$.
\end{fact}

\begin{lemma}
\label{lemma:disjoint-marking}
% Let the width of clauses in $F$ be $k=3$ (we can assume this without loss of generality).
Let $u\in T$ be a node at level greater than $m$. Then, at least one of the following must be true:
\begin{enumerate}
  \item
  $u$ has at most $2$ edges going out. 

  \item
  $u$ has one edge $e$ going out such that $|M(e)| \ge 2$.

  \item
  $u$ has one edge $e$ going out such that $Q_e\in X_D$ and $|M(e)| = 1$.
\end{enumerate}
\end{lemma}

\begin{proof}
Set variables that are part of $Q_u$ to $1$ and let $F'$ be the simplified $3$-CNF.
Say case $1$ does not happen. Then, the monotone clause $C$ used to develop edges out of $u$ has width $3$ and so, $C$ is also present in $F$. By \cref{fact:disjoint-intersect}, $u$ contains an edge $e$ such that $Q_e = x$ and $x\in X_D$. If $M(e) \ge 2$, then case $2$ is satisfied and if $M(e) = 1$, then case $3$ is satisfied.
\end{proof}


\subsection{An Analysis of $\ts$ for arbitrary $3$-CNF}

We extend \cref{path-pessim} for arbitrary $3$-CNFs taking into double markings into account.

\begin{lemma}
\label{lemma:double-marking-weight}
Let $T$ be a transversal tree and let $P = P_{ru}$ be a path starting from root $r$. Then $\sigma(P) \le (\frac{1}{\sqrt{3}})^{W^{+}(P) + W_{\ge 2}(P)}$.
\end{lemma}

\begin{proof}
For a marked edge $e \in P$, we define the contribution of $e$ as
\[
q_e := \prod_{v \in M(e)}\left(\frac{1}{|N_u(v)|+1}\right)^{1/|N_u(v)|}.
\]
By Fact \ref{fact:path-prob}, $\total(P) = \prod_{e : M(e) \ne \emptyset} q_e$. It is then sufficient to show that $q_e \le \lambda$ for every marked edge $e$. Note that $q_e$ can be written as $(\frac{1}{2})^a (\frac{1}{\sqrt{3}})^b$, for some non-negative integers $a$ and $b$ such that $a + b \ge |M(e)|$. This quantity is at most $\frac{1}{\sqrt{3}}$ if $|M(e)| = 1$ and is at most $\frac{1}{3}$ if $|M(e)| \ge 2$. Finally, we observe that $e$ contributes $1$ to $W(P)$ if $|M(e)| \ge 1$ and contributes $1$ to $W_{\ge 2}(P)$ if $M(e) \ge 2$.
\end{proof}



\begin{definition}[$NM(w, d, \openinteger)$]
For non-negative integers $w, d, \openinteger$, define $NM(w, d, \openinteger)$ to be the maximum of sum of survival probabilities of leaves over depth-$d$ transversal trees $T$ for $3$-CNFs. Moreover, for every root to leaf shoot $S$, $W^+(S) \ge w$ and $\openfunction(S) \le \openinteger$.
\end{definition}

% We not  basic facts regarding $NM(w, d, \open)$.

% \begin{fact}
% \label{fact:open-less-smaller}
%    If $\open\ge 1$, then $NM(w, d, \open) \ge NM(w, d, \open-1)$. 
% \end{fact}

% \begin{proof}
% By definition, transversal trees of weight $w$, depth $d$, and openness $\open-1$ are subset of transversal trees of weight $w$, depth $d$, and openness $\open$. The claim then follows.
% \end{proof}

% \begin{fact}
% \label{fact:weight-more-bigger}
% If $w\ge 1$, then $NM(w, d, \open) \le NM(w-1, d, \open)$. 
% \end{fact}


\subsubsection{Proving \cref{thm:main-result}}

We will show the following as our main lemma:

\begin{lemma}
\label{pess-nonmonotone-ntran}
Let $F$ be a $3$-CNF over $n$ variables with $\trn(F) = t \le \frac{n}{2}$. Then for any canonical transversal tree $T$ for $F$, it holds that 
\[
\total(T) \le
\begin{cases}
3^t & t \le \frac{n}{3}\\
3^{\frac{t}{3}}\times M'\left(3t - n, \frac{2t}{3}\right) & \textrm{ otherwise}
\end{cases}
\]
where $M'(w, d)$ is the same bound we obtained in \cref{sec:monotone}.
% Thus $R(n, t, 3) \le 3^{\frac{t}{3}}\times M(3t - n, \frac{2}{3}t)$, and in particular, $R(n, \frac{n}{2}, 3) \le 1.598^n$.
\end{lemma}


Using this, \cref{thm:main-result} follows from \cref{pess-nonmonotone-ntran} by using the exact same argument as \hyperref[proof:monotone-main-result]{in the Proof of monotone case of Theorem 1}.


Our main lemma will make use of the following bounds on $NM(w, d, \openinteger)$:

\begin{lemma} 
\label{lem:NMwdo}
For all $0\le d\le n, 0\le w\le 3d, 0\le \openinteger\le d$, it holds that:
\[
NM(w,d,\openinteger) \le
\begin{cases}
(2+\lambda)^{\openinteger}(2+\lambda^2)^{d-\openinteger}
 & 0\le w \le d\\
(2+\lambda)^{\openinteger - (w-d)}(1+2\lambda)^{w-d}(2+\lambda^2)^{d-\openinteger}
 & d\le w \le d + \openinteger\\
(1+2\lambda)^{\openinteger}(2+\lambda^2)^{2d-w}(1+\lambda+\lambda^2)^{w-d-\openinteger}
& d+\openinteger\le w\le 2d\\
(1+2\lambda)^{\openinteger - (w-2d)}(3\lambda)^{w-2d}(1+\lambda+\lambda^2)^{d-\openinteger}
& 2d\le w\le 2d+\openinteger\\
(3\lambda)^{\openinteger}(1+\lambda+\lambda^2)^{3d-w}(2\lambda + \lambda^2)^{w-2d-\openinteger}
& 2d+\openinteger\le w\le 3d\\
\end{cases}
\]
Moreover, for $\openinteger \ge d$:
\[
   NM(w, d, \openinteger) \le M'(w, d)
\]
where $M'(w, d)$ is from $\cref{sec:monotone}$.
\end{lemma}


Using this, we now prove our main lemma, which yields \cref{thm:main-result} as desired.

\begin{proof}[\textbf{Proof of \cref{pess-nonmonotone-ntran} assuming \cref{lem:NMwdo}}]
If $t \le \frac{n}{3}$, then observe that $T$ has at most $3^t$ leaves and we trivially bound $\total(T)\le 3^t$.

For $t \ge \frac{n}{3}$, we proceed by considering the maximal set of disjoint monotone width $3$ clauses in $F$ used to develop the first $m$ levels of $T$.
For every node $u\in T$ at level $m$, let the subtree rooted at $u$ be $T_u$.
We can bound $\total(T_u) \le NM(w, d, \openinteger)$ where $d = t-m, w = 3t-n, \openinteger = 2m$. Hence, $\total(T) \le 3^m NM(w, d, \openinteger)$.

If $m \ge \frac{t}{3}$, then $\openinteger = 2m \ge t - m = d$. Applying \cref{lem:NMwdo}, we infer that
\begin{align*}
\total(T) & \le 3^m M'(3t-n, t-m)\\
& = 3^{t/3} \left(3^{m-t/3} M'\left(3t-n, \frac{2t}{3} - \left(m - \frac{t}{3}\right) \right)\right)\\
& \le 3^{t/3} M'\left(3t-n, \frac{2t}{3}\right)\\
\end{align*}
and we infer the claim.

So, we assume that $m \le \frac{t}{3}$ and try to find the value of $m$ which will maximize $\total(T)$. Notice that in this case, $\openinteger \le d$ and so, we can't directly reduce to the case of $M'(w, d)$. 
As $t\le \frac{n}{2}$, it must be that $w = 3t-n \le t\le t+m \le d + \openinteger$.
This implies $w\le d+\openinteger$.
We now take two cases based on value of $w$ and apply \cref{lem:NMwdo} for the case of $\openinteger \le d$.

\begin{enumerate}[label=\textbf{Case \arabic*:}, itemindent=*, leftmargin=0em]

\item $0\le w\le d$.
In this case, we see that:
\begin{align*}
\total(T)   
   & \le 3^m NM\left(3t-n, t - m, 2m\right)\\
   & \le 3^m (2+\lambda)^{2m}(2+\lambda^2)^{t-3m}\\
   & = (2+\lambda^2)^t\left(\frac{(3)(2+\lambda)^2}{(2+\lambda^2)^3}\right)^m\\
\end{align*}
Here, the fraction has value $> 1$ and so is maximized when $m$ is maximized, i.e., when $m = \frac{t}{3}$ in which case:
\begin{align*}
\total(T)
& \le 3^{t/3} (2+\lambda)^{2t/3}\\
& = 3^{t/3} M'\left(3t-n, \frac{2t}{3}\right)
\end{align*}
Here the last equality follows by considering the case of $0\le w\le d$ for $M'(w, d)$.

\item $d\le w\le d + \openinteger$.
In this case, we see that:
\begin{align*}
\total(T)
& \le 3^m NM\left(3t-n, t - m, 2m\right)\\
& \le 3^m (2+\lambda)^{n-2t+m}(1+2\lambda)^{2t-n+m}(2+\lambda^2)^{t-3m}\\
& = \left(\frac{2+\lambda}{1+2\lambda}\right)^{n-2t}(2+\lambda^2)^t\left(\frac{(3)(2+\lambda)(1+2\lambda)}{(2+\lambda^2)^3}\right)^m\\
\end{align*}
Here, the rightmost fraction is $> 1$ and so is maximized when $m$ is maximized, i.e., when $m = \frac{t}{3}$ in which case:
\begin{align*}
\total(T)
& \le 3^{t/3} (2+\lambda)^{n - 5t/3} (1+2\lambda)^{7t/3 - n}\\
& = 3^{t/3} M'\left(3t-n, \frac{2t}{3}\right)
\end{align*}
Here the last equality follows by considering the case of $d\le w\le 2d$ for $M'(w, d)$ (we can do this as $\openinteger \le d$ and hence, $w\le 2d$).
\end{enumerate}

Thus, in either case, we exactly recover the monotone bound as desired.
\end{proof}








\subsubsection{Upper bounds on $NM(w, d, \openinteger)$}

As done in monotone analysis, we let $\lambda \eqdef \frac{1}{\sqrt{3}}$. 
Our goal is to prove \cref{lem:NMwdo}.
We introduce a recurrence relation $L(w, d, \openinteger)$ that we argue will upper bound $NM(w, d, \openinteger)$.

\begin{definition}[$L(w, d, \openinteger)$]
\label{def:Lwdo}
We define $L(w, d, \openinteger): \N^3\rightarrow \R$ recursively as follows:
\[
L(w, 0, \openinteger) = 
\begin{cases}
1 & w \le 0\\
0 & w > 0\\
\end{cases}
\]
For $w\le 3d, 0\le d\le n$, and $\openinteger \ge 1$, define $L(w, d, \openinteger)$ as:
\begin{align*}
L(w, d, \openinteger) = 
\max\{
& (2+\lambda) L(w-1, d-1, \openinteger-1),\\
& (1+2\lambda) L(w-2, d-1, \openinteger-1),\\
& 3\lambda L(w-3, d-1, \openinteger-1)
\}\\
\end{align*}
For $w\le 3d, 0\le d\le n$, and $\openinteger = 0$, define $L(w, d, 0)$ as:
\begin{align*}
L(w, d, 0) =
\max\{
& (2+\lambda^2) L(w-1, d-1, 0),\\
& (1+\lambda+\lambda^2) L(w-2, d-1, 0),\\
& (2\lambda + \lambda^2) L(w-3, d-1, 0)
\}\\
\end{align*}
\end{definition}

We claim that $L(w, d, \openinteger)$ gives a good bound on $M(w, d, \openinteger)$.

\begin{prop} 
\label{claim:NMwd-recurrence}
For all $0\le d\le n, 0\le w\le 3d, 0\le \openinteger$, it holds that:
$NM(w, d, \openinteger) \le L(w, d, \openinteger)$
\end{prop}

We will show the following bound on $L(w, d, \openinteger)$.

\begin{lemma} 
\label{lem:Lwdo}
For all $0\le d\le n, 0\le w\le 3d, 0\le \openinteger\le d$, it holds that:
\[
L(w,d,\openinteger) \le
\begin{cases}
(2+\lambda)^{\openinteger}(2+\lambda^2)^{d-\openinteger}
 & 0\le w \le d\\
(2+\lambda)^{\openinteger - (w-d)}(1+2\lambda)^{w-d}(2+\lambda^2)^{d-\openinteger}
 & d\le w \le d + \openinteger\\
(1+2\lambda)^{\openinteger}(2+\lambda^2)^{2d-w}(1+\lambda+\lambda^2)^{w-d-\openinteger}
& d+\openinteger\le w\le 2d\\
(1+2\lambda)^{\openinteger - (w-2d)}(3\lambda)^{w-2d}(1+\lambda+\lambda^2)^{d-\openinteger}
& 2d\le w\le 2d+\openinteger\\
(3\lambda)^{\openinteger}(1+\lambda+\lambda^2)^{3d-w}(2\lambda + \lambda^2)^{w-2d-\openinteger}
& 2d+\openinteger\le w\le 3d\\
\end{cases}
\]
Moreover, for $\openinteger \ge d$:
\[
   L(w, d, \openinteger) \le M'(w, d)
\]
where $M'(w, d)$ is from $\cref{sec:monotone}$.
\end{lemma}

Combining \cref{claim:NMwd-recurrence} and \cref{lem:Lwdo}, \cref{lem:NMwdo} trivially follows.


\subsubsection{Proving $NM(w, d, \openinteger)\le L(w, d, \openinteger)$}


Using \cref{lemma:disjoint-marking} and \cref{lemma:double-marking-weight}, we come up with a recurrence for $NM(w, d, \openinteger)$ and show it's bounded by $L(w, d, \openinteger)$, proving \cref{claim:NMwd-recurrence}.

\begin{proof}[\textbf{Proof of \cref{claim:NMwd-recurrence}}]
Recall that in the canonical ordering, all width $3$ monotone clauses appear first and remaining clauses appear later. After exhausting the width $3$ monotone clauses, the remaining clauses that we develop in the transversal tree have width at most $2$. Towards this, for non-negative integers $w, d$: let $M_2(w, d)$ be the maximum sum of survival probabilities over all transversal trees for $2$-CNFs where every root to leaf path has uniform weight at least $w$. Recall that uniform weight is defined with respect to $3$-CNFs and we continue using that definition.
We get various recurrences by considering cases on number of marked edges out of the root node ($0$ or $1$ or $2$) and by observing that some cases are dominated by others (such as various cases of width $1$ clauses).
The remaining recurrences that are not dominated by any other recurrence are the following:
\begin{align*}
M_2(w, d) \le 
\max\{
& (2)M_2(w-1, d-1),\\
& (1+\lambda)M_2(w-2, d-1),\\
& 2\lambda M_2(w-3, d-1)
\}\\
\end{align*}
By induction, we infer that $M_2(w, d) \le L(w, d, 0)$.

We now develop a recurrence for $NM(w, d, 0)$.
Recall that in canonical ordering, either we exhaust all width $3$ monotone clause and reach $M_2(w, d)$, or we develop width $3$ monotone clause.
Observe that \cref{lemma:disjoint-marking} guarantees that every node in such a tree must have an edge $e$ coming out of it such that $|M(e)| \ge 2$. Taking cases on the number of marked edges coming out of the root node and whether root node has $3$ or at most $2$ edges coming out, we get many recurrences.
However certain recurrences are dominated by others and the remaining recurrences that are not dominated by any othe recurrence are as follows:
\begin{align*}
NM(w, d, 0) \le 
\max\{
& (2 + \lambda^2)NM(w-1, d-1, 0),\\
& (1+\lambda+\lambda^2)NM(w-2, d-1, 0),\\
& (2\lambda+\lambda^2) NM(w-3, d-1, 0),\\
& M_2(w, d)
\}\\
\end{align*}
By induction, we again infer that $NM(w, d, 0) \le L(w, d, 0)$.

We now develop a recurrence for $NM(w, d, \openinteger)$.
We again take advantange of the fact that in canonical ordering, either we exhaust all width $3$ monotone clause and reach $M_2(w, d)$, or we develop width $3$ monotone clause. Moreover, amongst width $3$ clauses, canonical ordering causes either $\openinteger$ to decrease by at least $1$ or we exhaust such clauses and all remaining clauses have the property that a node developed using such a clause will have an outgoing edge $e$ such that $|M(e)| \ge 2$. 

We get many recurrences for $NM(w, d, \openinteger)$ by considering cases on number of marked edges ($1$ or $2$ or $3$) out of the root node, number of marked edges that cause $\openfunction$ to decrease ($1$ or $2$ or $3$), various combinations of number of double marked edges ($1$ or $2$ or $3$), whether the root node has at most $2$ edges coming out, and whether the root node has no edges that cause $\openfunction$ to decrease by at least $1$. Notice that if the root node has no edges that cause $\openfunction$ to decrease by at least $1$, then by clause ordering, no other width $3$ clause can cause $\openfunction$ to decrease and hence, we are in case $NM(w, d, 0)$. Lastly, we observe that certain recurrences are dominated by others. The remaining recurrences that are not dominated by any other recurrences are the following:
\begin{align*}
NM(w, d, \openinteger) \le 
\max\{
& (2+\lambda) NM(w-1, d-1, \openinteger-1),\\
& (1+2\lambda) NM(w-2, d-1, \openinteger-1),\\
& 3\lambda NM(w-3, d-1, \openinteger-1),\\
& NM(w, d, 0),\\
& M_2(w, d)
\}\\
\end{align*}
By induction, utilizing the fact that $L(w, d, \openinteger) \ge L(w, d, 0)$, we again infer that $NM(w, d, \openinteger) \le L(w, d, \openinteger)$ as desired.
\end{proof}





\subsubsection{Upper bound on $L(w, d, 0)$}

We first show \cref{lem:Lwdo} for the special case of $\openinteger = 0$:

\begin{lemma} 
\label{lem:Lwdo-0}
For all $0\le d\le n, 0\le w\le 3d$, it holds that:
\[
L(w,d,0) \le
\begin{cases}
(2+\lambda^2)^d & 0\le w \le d\\
(2+\lambda^2)^{2d-w}(1+\lambda+\lambda^2)^{w-d} & d\le w\le 2d\\
(1+\lambda+\lambda^2)^{3d-w}(2\lambda+\lambda^2)^{w-2d} & 2d\le w\le 3d\\
\end{cases}
\]
\end{lemma}

\begin{proof}
% [\textbf{Proof of \cref{lem:Lwdo-0}}]

Let $G_1, G_2, G_3, G: \N^2\rightarrow \R$ be defined as:

\begin{eqnarray*}
G_1(w, d) & = & (2+\lambda^2)^d\\
G_2(w, d) & = & (2+\lambda^2)^{2d-w}(1+\lambda+\lambda^2)^{w-d}\\
G_3(w, d) & = & (1+\lambda+\lambda^2)^{3d-w}(2\lambda+\lambda^2)^{w-2d}\\
G(w, d) & = & \min\{G_1(w, d), G_2(w, d), G_3(w, d)\}\\
\end{eqnarray*}

For $1\le i\le 3$, define $P_i$ to be the set of pairs $(w, d)$ such that $d \ge 0$ and $w\in [(i-1)d, id]$. We will show the following two propositions:

\begin{proposition}
\label{prop:G}
For all $1\le i\le 3$, and all $(w, d)\in P_i: G(w, d) = G_i(w, d)$.
\end{proposition}

\begin{proposition}
\label{prop:L-sol-open-0}
For all $1\le i\le 3$ and all $(w, d)\in P_i: L(w, d, 0) \le G(w, d)$.
\end{proposition}

We observe that \cref{prop:G} and \cref{prop:L-sol-open-0} together imply our claim.

\begin{proof}[\textbf{Proof of \cref{prop:G}}]
The result follows immediately from the following claims:
\begin{claim}
\label{claim:G-sol-1}
$G_1(w,d) \leq G_2(w,d)$ if and only if $w \leq d$, with equality when $w=d$.
\end{claim}

\begin{claim}
\label{claim:G-sol-2}
$G_2(w,d) \leq G_3(w,d)$ if and only if $w \leq 2d$, with equality when $w=2d$.
\end{claim}

\cref{claim:G-sol-1} holds because:

\begin{equation*}
\frac{G_1(w, d)}{G_2(w, d)} = \left(\frac{2+\lambda^2}{1+\lambda+\lambda^2}\right)^{w-d}
\end{equation*}

which is greater than 1 if and only if $w>d$.

\cref{claim:G-sol-2} holds because:

\begin{equation*}
\frac{G_2(w, d)}{G_3(w, d)} = \left(\frac{(1+\lambda+\lambda^2)^2}{(2\lambda + \lambda^2)(2+\lambda^2)}\right)^{w-2d}
\end{equation*}
which is greater than 1 if and only if $w>2d$.
\end{proof}


\begin{proof}[\textbf{Proof of \cref{prop:L-sol-open-0}}]
We consider cases on value of $w$ and in every case, induct on $d$ and apply \cref{def:Lwdo} to infer the claim.


\begin{enumerate}[label=\textbf{Case \arabic*:}, itemindent=*, leftmargin=0em]

\item Assume $(w, d)\in P_1$.
\begin{eqnarray*}
L(w,d,0) & \le & \max \{  (2+\lambda^2) G(w-1,d-1,0),\\
          &   &          (1+\lambda+\lambda^2) G(w-2,d-1,0),\\
          &   &          (2\lambda+\lambda^2) G(w-3,d-1,0) \}\\
& \le  & \max\{(2+\lambda^2)G_1(w-1,d-1),(1+\lambda+\lambda^2)G_1(w-2,d-1),(2\lambda+\lambda^2)G_1(w-3,d-1)\}\\
& = & G_1(w,d) \max\{1, (1+\lambda+\lambda^2)/(2+\lambda^2), (2\lambda+\lambda^2)/(2+\lambda^2)\}\\
& = & G_1(w,d)\\
& = & G(w, d)
\end{eqnarray*}
The last equality follows by applying \cref{prop:G} for the case $(w, d)\in P_1$.

\item Assume $(w, d)\in P_2$.
\begin{eqnarray*}
L(w,d,0) & \le & \max\{(2+\lambda^2) G(w-1,d-1,0),\\
          &   &        (1+\lambda+\lambda^2) G(w-2,d-1,0),\\
          &   &       (2\lambda+\lambda^2) G(w-3,d-1,0) \} \\
& \le & \max \{(2+\lambda^2)G_2(w-1,d-1),(1+\lambda+\lambda^2)G_2(w-2,d-1),(2\lambda+\lambda^2)G_2(w-3,d-1)\}\\
& = & G_2(w,d) \max\{1, 1, (2\lambda+\lambda^2)(2+\lambda^2)/(1+\lambda+\lambda^2)^2\}\\
& = & G_2(w,d)\\
& = & G(w, d)
\end{eqnarray*}
The last equality follows by applying \cref{prop:G} for the case $(w, d)\in P_2$.


\item Assume $(w, d)\in P_3$.
\begin{eqnarray*}
L(w,d,0) & \le & \max\{(2+\lambda^2) G(w-1,d-1,0),\\
         &   &        (1+\lambda+\lambda^2) G(w-2,d-1,0),\\
         &   &        (2\lambda+\lambda^2) G(w-3,d-1,0) \}\\
& \le & \max \{(2+\lambda^2)G_3(w-1,d-1),(1+\lambda+\lambda^2)G_3(w-2,d-1),(2\lambda+\lambda^2)G_3(w-3,d-1)\}\\
& = & G_3(w,d) \max\{(2+\lambda^2)(2\lambda+\lambda^2)/(1+\lambda+\lambda^2)^2, 1, 1\}\\
& = & G_3(w,d)\\
& = & G(w,d)\\
\end{eqnarray*}
The last equality follows by applying \cref{prop:G} for the case $(w, d)\in P_3$.


\end{enumerate}


\end{proof}

\end{proof}





















\subsubsection{Upper bound on $L(w, d, \openinteger)$}

We are finally ready to give general bounds on $L(w, d, \openinteger)$:

\begin{proof}[\textbf{Proof of \cref{lem:Lwdo}}]

Notice that if $\openinteger \ge d$, then if we try and unravel the recurrence, no path can lead to the case $\openinteger = 0, d > 0$. Hence, $\openinteger$ plays no role in restricting the recurrence and $L(w, d, \openinteger)$ follows the same recurrence as $M'(w, d)$, yielding the claim.

For $\openinteger \le d$, we proceed by first defining $H_1, H_2, H_3, H_4, H_5, H: \N^3\rightarrow \R$ as follows:

\begin{eqnarray*}
H_1(w, d, \openinteger) & = & (2+\lambda)^{\openinteger}(2+\lambda^2)^{d-\openinteger}\\
H_2(w, d, \openinteger) & = & (2+\lambda)^{\openinteger - (w-d)}(1+2\lambda)^{w-d}(2+\lambda^2)^{d-\openinteger}\\
H_3(w, d, \openinteger) & = & (1+2\lambda)^{\openinteger}(2+\lambda^2)^{2d-w}(1+\lambda+\lambda^2)^{w-d-\openinteger}\\
H_4(w, d, \openinteger) & = & (1+2\lambda)^{\openinteger - (w-2d)}(3\lambda)^{w-2d}(1+\lambda+\lambda^2)^{d-\openinteger}\\
H_5(w, d, \openinteger) & = & (3\lambda)^{\openinteger}(1+\lambda+\lambda^2)^{3d-w}(2\lambda + \lambda^2)^{w-2d-\openinteger}\\
H(w, d, \openinteger) & = & \min\{H_2(w, d, \openinteger), H_3(w, d, \openinteger), H_4(w, d, \openinteger), H_5(w, d, \openinteger)\}\\
\end{eqnarray*}

For $1\le i\le 5$, define $Q_i\subset \N^3$ as follows:

\begin{align*}
Q_1 & = \{(w, d, \openinteger)\in \N^3: 0 \le w \le d + \openinteger\}\\
Q_2 & = \{(w, d, \openinteger)\in \N^3: d \le w \le d + \openinteger\}\\
Q_3 & = \{(w, d, \openinteger)\in \N^3: d + \openinteger \le w \le 2d\}\\
Q_4 & = \{(w, d, \openinteger)\in \N^3: 2d \le w \le 2d + \openinteger\}\\
Q_5 & = \{(w, d, \openinteger)\in \N^3: 2d + \openinteger \le w \le 3d\}\\
\end{align*}

 We will show the following propositions that together imply our claim:

\begin{proposition}
\label{prop:H}
For all $1\le i\le 5$, and all $(w, d, \openinteger)\in Q_i: H(w, d, \openinteger) = H_i(w, d, \openinteger)$.
\end{proposition}


\begin{proposition}
\label{prop:L-sol-open}
For all $1\le i\le 5$ and all $(w, d, \openinteger)\in Q_i: L(w, d, \openinteger) \le H(w, d, \openinteger)$.
\end{proposition}

We will in fact use \cref{prop:H} in the proof of \cref{prop:L-sol-open}. Hence, we prove the former first:


\begin{proof}[\textbf{Proof of \cref{prop:H}}]
The result follows immediately from the following claims:

\begin{claim}
\label{claim:H-sol-1}
$H_1(w, d, \openinteger) \leq H_2(w, d, \openinteger)$ if and only if $w \leq d$, with equality when $w = d$.
\end{claim}

\begin{claim}
\label{claim:H-sol-2}
$H_2(w, d, \openinteger) \leq H_3(w, d, \openinteger)$ if and only if $w \leq d + \openinteger$, with equality when $w = d + \openinteger$.
\end{claim}

\begin{claim}
\label{claim:H-sol-3}
$H_3(w, d, \openinteger) \leq H_4(w, d, \openinteger)$ if and only if $w \leq 2d$, with equality when $w = 2d$.
\end{claim}

\begin{claim}
\label{claim:H-sol-4}
$H_4(w, d, \openinteger) \leq H_5(w, d, \openinteger)$ if and only if $w \leq 2d+\openinteger$, with equality when $w = 2d+\openinteger$.
\end{claim}

\cref{claim:H-sol-1} holds because: 

\begin{equation*}
\frac{H_1(w, d, \openinteger)}{H_2(w, d, \openinteger)} = \left(\frac{2+\lambda}{1+2\lambda}\right)^{w-d}
\end{equation*}

which is greater than 1 if and only if $w > d$.

\cref{claim:H-sol-2} holds because: 

\begin{equation*}
\frac{H_2(w, d, \openinteger)}{H_3(w, d, \openinteger)} = \left(\frac{(1+2\lambda)(2+\lambda^2)}{(2+\lambda)(1+\lambda+\lambda^2)}\right)^{w-d-\openinteger}
\end{equation*}

which is greater than 1 if and only if $w > d + \openinteger$.

\cref{claim:H-sol-3} holds because:

\begin{equation*}
\frac{H_3(w, d, \openinteger)}{H_4(w, d, \openinteger)} = \left(\frac{(1+2\lambda)(1+\lambda+\lambda^2)}{(2+\lambda^2)(3\lambda)}\right)^{w-2d}
\end{equation*}

which is greater than 1 if and only if $w>2d$.

\cref{claim:H-sol-4} holds because:

\begin{equation*}
\frac{H_4(w, d, \openinteger)}{H_5(w, d, \openinteger)} =\left(\frac{(3\lambda)(1+\lambda+\lambda^2)}{(1+2\lambda)(2\lambda+\lambda^2)}\right)^{w-2d-\openinteger}
\end{equation*}

which is greater than 1 if and only if $w>2d+\openinteger$.
\end{proof}

We prove our final proposition:

\begin{proof}[\textbf{Proof of \cref{prop:L-sol-open}}]
We observe that for $\openinteger = 0$, our claim follows from \cref{lem:Lwdo-0}.
We use this fact in the inductive argument below and only consider cases where $\openinteger \ge 1$. We consider cases on value of $w$ and in every case, induct on $d + \openinteger$ and apply \cref{def:Lwdo} to infer the claim.  

\begin{enumerate}[label=\textbf{Case \arabic*:}, itemindent=*, leftmargin=0em]

\item Assume $(w, d, \openinteger)\in Q_1$ and $\openinteger \ge 1$.
\begin{eqnarray*}
L(w, d, \openinteger) & = & \max\{(2+\lambda)H(w-1, d-1, \openinteger-1), (1+2\lambda)H(w-2, d-1, \openinteger-1), \\
& & (3\lambda)H(w-3, d-1, \openinteger-1)\}\\
& \le & \max\{(2+\lambda)H_1(w-1, d-1, \openinteger-1), (1+2\lambda)H_1(w-2, d-1, \openinteger-1), \\
& & (3\lambda)H_1(w-3, d-1, \openinteger-1)\}\\
& = & H_1(w, d, \openinteger)\max\left\{1, \frac{1+2\lambda}{2+\lambda}, \frac{3\lambda}{2+\lambda}\right\} \\
& = & H_1(w, d, \openinteger)\\
& = & H(w, d, \openinteger)
\end{eqnarray*}
The last equality follows by applying \cref{prop:H} for the case $(w, d,\openinteger)\in Q_1$.

\item Assume $(w, d, \openinteger)\in Q_2$ and $\openinteger \ge 1$.
\begin{eqnarray*}
L(w, d, \openinteger) & \leq & \max\{(2+\lambda)L(w-1, d-1, \openinteger-1), (1+2\lambda)L(w-2, d-1, \openinteger-1), \\
& & (3\lambda)L(w-3, d-1, \openinteger-1)\}\\
& \le & \max\{(2+\lambda)H_2(w-1, d-1, \openinteger-1), (1+2\lambda)H_2(w-2, d-1, \openinteger-1), \\
& & (3\lambda)H_2(w-3, d-1, \openinteger-1)\}\\
& = & H_2(w, d, \openinteger)\max\left\{1, 1, \frac{(3\lambda)(2+\lambda)}{(1+2\lambda)^2} \right\}\\
& = & H_2(w, d, \openinteger)\\
& = & H(w, d, \openinteger)\\
\end{eqnarray*}
The last equality follows by applying \cref{prop:H} for the case $(w, d,\openinteger)\in Q_2$.

\item Assume $(w, d, \openinteger)\in Q_3$ and $\openinteger \ge 1$.
\begin{eqnarray*}
L(w, d, \openinteger) & \leq & \max\{(2+\lambda)L(w-1, d-1, \openinteger-1), (1+2\lambda)L(w-2, d-1, \openinteger-1), \\
& & (3\lambda)L(w-3, d-1, \openinteger-1)\}\\
& \le & \max\{(2+\lambda)H_3(w-1, d-1, \openinteger-1), (1+2\lambda)H_3(w-2, d-1, \openinteger-1), \\
& & (3\lambda)H_3(w-3, d-1, \openinteger-1)\}\\
& = & H_3(w, d, \openinteger)\max\left\{\frac{(2+\lambda)(1+\lambda+\lambda^2)}{(1+2\lambda)(2+\lambda^2)}, 1, \frac{(2+\lambda^2)(3\lambda)}{(1+2\lambda)(1+\lambda+\lambda^2)} \right\}\\
& = & H_3(w, d, \openinteger)\\
& = & H(w, d, \openinteger)\\
\end{eqnarray*}
The last equality follows by applying \cref{prop:H} for the case $(w, d,\openinteger)\in Q_3$.

\item Assume $(w, d, \openinteger)\in Q_4$ and $\openinteger \ge 1$.
\begin{eqnarray*}
L(w, d, \openinteger) & \leq & \max\{(2+\lambda)L(w-1, d-1, \openinteger-1), (1+2\lambda)L(w-2, d-1, \openinteger-1), \\
& & (3\lambda)L(w-3, d-1, \openinteger-1)\}\\
& \le & \max\{(2+\lambda)H_4(w-1, d-1, \openinteger-1), (1+2\lambda)H_4(w-2, d-1, \openinteger-1), \\
& & (3\lambda)H_4(w-3, d-1, \openinteger-1)\}\\
& = & H_4(w, d, \openinteger)\max\left\{\frac{(3\lambda)(2+\lambda)}{(1+2\lambda)^2}, 1, 1 \right\}\\
& = & H_4(w, d, \openinteger)\\
& = & H(w, d, \openinteger)
\end{eqnarray*}
The last equality follows by applying \cref{prop:H} for the case $(w, d,\openinteger)\in Q_4$.

\item Assume $(w, d, \openinteger)\in Q_5$ and $\openinteger \ge 1$.
\begin{eqnarray*}
L(w, d, \openinteger) & \leq & \max\{(2+\lambda)L(w-1, d-1, \openinteger-1), (1+2\lambda)L(w-2, d-1, \openinteger-1), \\
& & (3\lambda)L(w-3, d-1, \openinteger-1)\}\\
& \le & \max\{(2+\lambda)H_5(w-1, d-1, \openinteger-1), (1+2\lambda)H_5(w-2, d-1, \openinteger-1), \\
& & (3\lambda)H_5(w-3, d-1, \openinteger-1)\}\\
& = & H_5(w, d, \openinteger)\max\left\{\frac{(2+\lambda)(2\lambda+\lambda^2)^2}{(3\lambda)(1+\lambda+\lambda^2)^2},
                             \frac{(1+2\lambda)(2\lambda+\lambda^2)}{(3\lambda)(1+\lambda+\lambda^2)}, 1 \right\}\\
& = & H_5(w, d, \openinteger)\\
& = & H(w, d, \openinteger)
\end{eqnarray*}
The last equality follows by applying \cref{prop:H} for the case $(w, d,\openinteger)\in Q_5$.

\end{enumerate}

\end{proof}

\end{proof}

























% \subsection{Getting the actual bound}

% We now recover the same bound as we got in monotone analysis earlier.

% \begin{lemma}
% \label{pess-nonmonotone-ntran}
% Let $F$ be a $3$-CNF over $n$ variables with $\trn(F) = t \le \frac{n}{2}$. Then for any canonical transversal tree $T$ for $F$, it holds that 
% \[
% \total(T) \le
% \begin{cases}
% 3^t & t \le \frac{n}{3}\\
% 3^{\frac{t}{3}}\times M(3t - n, \frac{2}{3}t) & \textrm{ otherwise}
% \end{cases}
% \]
% % Thus $R(n, t, 3) \le 3^{\frac{t}{3}}\times M(3t - n, \frac{2}{3}t)$, and in particular, $R(n, \frac{n}{2}, 3) \le 1.598^n$.
% \end{lemma}

% \begin{proof}
% We proceed by considering any maximal set of disjoint monotone width $3$ clauses in $F$ and let $m$ be the size of this set. If $m > \frac{t}{3}$, then we exactly recover the monotone bound.\\

% So, we assume that $m \le \frac{t}{3}$ and try to find the value of $m$ which will maximize $\total(T)$.
% For every node $u$ at level $m$, $d = t-m, w = \max(3t-n, d) = \max(3t-n, t-m),  \open = 2m$.
% As $t\le \frac{n}{2}$, it must be that $3t-n \le t\le t+m \le d + \open$.
% This implies $w\le d+\open$.
% We now take two cases based on value of $w$ and apply \cref{lem:Rwdo}:
% \begin{enumerate}[label=\textbf{Case \arabic*:}, itemindent=*, leftmargin=0em]
%    \item $w = d$. Then,
%    \begin{align*}
%    \total(T)   
%       & \le 3^m NM\left(d, d, \open\right)\\
%       & \le 3^m R\left(d, d, \open\right)\\
%       & = 3^m (2+\lambda)^{2m}(2+\lambda^2)^{t-3m}\\
%       & = (2+\lambda^2)^t\left(\frac{(3)(2+\lambda)^2}{(2+\lambda^2)^3}\right)^m\\
%    \end{align*}
%    Here, the fraction has value $> 1$ and so is maximized when $m$ is maximized, i.e., when $m = \frac{t}{3}$ in which case we recover the monotone bound.

%    \item $w > d$. Then,
%    \begin{align*}
%    \total(T)
%       & \le 3^m NM\left(w, d, \open\right)\\
%       & \le 3^m R\left(w, d, \open\right)\\
%       & = 3^m (2+\lambda)^{n-2t+m}(1+2\lambda)^{2t-n+m}(2+\lambda^2)^{t-3m}\\
%       & = \left(\frac{2+\lambda}{1+2\lambda}\right)^{n-2t}(2+\lambda^2)^t\left(\frac{(3)(2+\lambda)(1+2\lambda)}{(2+\lambda^2)^3}\right)^m\\
%    \end{align*}
%    Here, the fraction is $> 1$ and so is maximized when $m$ is maximized, i.e., when $m = \frac{t}{3}$ in which case we again recover the monotone bound.
% \end{enumerate}
% Hence, in either case, we exactly recover the monotone bound as desired.
% \end{proof}

% Now, \cref{thm:main-result} follows from \cref{pess-nonmonotone-ntran} by using the exact same argument as \hyperref[proof:monotone-main-result]{in the Proof of monotone case of Theorem 1}.























